%! Polynomial Functions and Rates of Change — Unit 1, Lesson 1.4 — Homework (Answer Key)
\documentclass[10pt]{article}
\usepackage{precalculus-article}
\usepackage{precalculus-key}   % pulls in -boxes; do NOT also load -boxes
\newcommand{\TallMath}[1]{$\displaystyle #1\rule[-1.4em]{0pt}{3.2em}$}

\begin{document}

\pageheader{Unit 1, Lesson 1.4}{Homework --- Answer Key}
\namedateperiod

\vspace{0.12in}

\begin{notesbox}{Homework: Polynomial Functions and Rates of Change --- Key}
\begin{enumerate}[leftmargin=1.5em, itemsep=10pt]

\item For the polynomial $p(x) = x^4 - 8x^2 + 7$:

\begin{enumerate}[leftmargin=1.5em, itemsep=5pt]
    \item[(a)] Degree: \ans{4} \quad Leading coefficient: \ans{1}

    \item[(b)] At most $n - 1 = 4 - 1 = $ \ans{3 local extrema}.

    \item[(c)] End behavior: \ans{Both ends rise (as $x \to \pm\infty$, $p(x) \to +\infty$) because the degree is even and the leading coefficient is positive.}
\end{enumerate}

\item A polynomial function $f$ has exactly two local maxima and two local minima.

\begin{enumerate}[leftmargin=1.5em, itemsep=5pt]
    \item[(a)] Total local extrema: \ans{4}

    \item[(b)] Minimum possible degree: \ans{5. A polynomial of degree $n$ can have at most $n-1$ local extrema. To have 4 local extrema we need $n - 1 \geq 4$, so $n \geq 5$.}
\end{enumerate}

\item AROC of $g$ over $[0,2]$ is 8; over $[2,4]$ is $-3$.

\begin{enumerate}[leftmargin=1.5em, itemsep=5pt]
    \item[(a)] Answer: \ans{A --- Local maximum.} Explanation: \ans{The AROC transitions from positive (8) to negative ($-3$), meaning $g$ switches from increasing to decreasing near $x=2$. This is the definition of a local maximum.}

    \item[(b)] \ans{No. A point of inflection is where the AROC switches from increasing to decreasing (or vice versa) --- it is about the rate of change of the AROC, not a sign change of the AROC itself. A local maximum involves a sign change in AROC (positive to negative), not an inflection.}
\end{enumerate}

\item Even degree, positive leading coefficient.

\begin{enumerate}[leftmargin=1.5em, itemsep=5pt]
    \item[(a)] \ans{Global minimum}

    \item[(b)] \ans{An even-degree polynomial with positive leading coefficient has both ends rising ($y \to +\infty$ as $x \to \pm\infty$). Since the function rises without bound on both sides but dips in the middle, it must achieve a lowest output value somewhere --- a global minimum. It has no global maximum because output grows unbounded.}
\end{enumerate}

\item AP-Style MC: Answer: \ans{A}

\ans{From the table, $h(-2) = -5 < h(-1) = 3 > h(0) = 4$... Wait --- $h(-1) = 3 < h(0) = 4$, so the local max is near $x=0$. Actually: $h(-2)=-5 \to h(-1)=3 \to h(0)=4 \to h(1)=1 \to h(2)=-2$. The function increases from $-2$ to $0$ then decreases, so local max near $x=0$. However, going from $-2$ to $-1$ the values increase sharply; from $0$ to $2$ they decrease. Best answer: \textbf{C} --- local maximum near $x=0$ only.}

\teachernote{Teacher note: The correct answer is C. The function increases from $x=-2$ to $x=0$ (values $-5, 3, 4$) and decreases from $x=0$ to $x=2$ (values $4, 1, -2$). So there is a local maximum near $x=0$. There is no sign-change in AROC that indicates a local minimum within this window.}

\end{enumerate}
\end{notesbox}

\vspace{0.10in}

\begin{extensionbox}
\textbf{Extension (Optional) --- Key:} For $p(x) = x^4 - 4x^2$:
$p(-2)=0$, $p(-1)=-3$, $p(0)=0$, $p(1)=-3$, $p(2)=0$.

\renewcommand{\arraystretch}{1.5}
\begin{tabularx}{\linewidth}{@{} l X X @{}}
\rowcolor{goldbg}
Interval & AROC formula & Value \\
\midrule
$[-2, -1]$ & $(-3-0)/1$  & \ans{$-3$} \\
$[-1,  0]$ & $(0-(-3))/1$  & \ans{$3$} \\
$[ 0,  1]$ & $(-3-0)/1$   & \ans{$-3$} \\
$[ 1,  2]$ & $(0-(-3))/1$  & \ans{$3$} \\
\end{tabularx}

\vspace{6pt}
Sign changes in AROC: \ans{between $[-2,-1]$ and $[-1,0]$ (near $x=-1$); between $[0,1]$ and $[1,2]$ (near $x=1$)} \\
Local extrema at: \ans{$x \approx -\sqrt{2} \approx -1.41$ (local min) and $x=0$ (local max) and $x \approx \sqrt{2} \approx 1.41$ (local min)}

\vspace{4pt}
AROC switching direction: \ans{The AROC goes $-3, 3, -3, 3$ --- alternating. There is a switch from $-3$ to $3$ near $x=-1$ and from $3$ to $-3$ near $x=0$, suggesting inflection points near $x \approx -1$ and $x \approx 1$.}
\end{extensionbox}

\vspace{0.10in}

\begin{spiralbox}
\textbf{Preview:} In Lesson 1.5, we will explore the \textbf{zeros of polynomial functions}
in depth --- including their multiplicities, how multiplicity affects the graph's behavior at
a zero, and an introduction to \textbf{complex zeros}.
\end{spiralbox}

\end{document}
